\problemname{Inflation}
A large part of the economy is influenced by the prevailing inflation, i.e., how quickly prices of
goods in society are rising. For example, the level of inflation affects how the Riksbank sets its
policy rate, which in turn impacts the stock market, mortgage rates, and so on.
Inflation is measured in percentage, indicating how much prices have increased compared to the same time last year.

In Sweden, the agency Statistics Sweden (SCB) measures the rate of inflation in society.
They calculate the inflation rate each month by measuring how the inflation rate \emph{changes every day}
during the month and then summing up all the daily increases.
You would like to use these figures to know if you should buy or sell all your stocks before SCB releases the figures officially.
Luckily, you have a friend who is willing to leak the figures to you.

One problem is that your friend has been quite inconsistent with his notes.
Sometimes he has recorded the daily change in the inflation rate in percentage points and sometimes in percent.
If the inflation rate was $20\%$ and the recorded increase was ``$50$'', you do not know if the new inflation
rate is $70\%$ (an increase of 50 percentage points) or $30\%$ (an increase of 50 percent).

At the very least, you can calculate an upper limit on the inflation rate, which might help your stock trading somewhat.
What is the highest possible inflation rate, given your friend's data?


\section*{Input}
The first line of input contains a decimal number $I$, the inflation rate in percent at the beginning of the month.
This is followed by a line with $30$ new decimal numbers, the changes in the inflation rate that your friend has
recorded during the month.

All numbers are between $0$ and $1000$ and have at most $2$ decimal places.

\section*{Output}
Print a decimal number: the highest possible inflation rate.
Your answer is considered correct if it has a relative or absolute error of at most $10^{-6}$.

\section*{Points}
Your solution will be tested on several test case groups.
To get the points for a group, it must pass all the test cases in the group.

\noindent
\begin{tabular}{| l | l | p{12cm} |}
  \hline
  \textbf{Group} & \textbf{Point value} & \textbf{Constraints} \\ \hline
  $1$    & $55$      & Your friend has written down all increases in percentage points. \\ \hline
  $2$    & $45$      & No additional constraints. \\ \hline
\end{tabular}

%\section*{Samples}
